\documentclass[report.tex]{subfiles}

\begin{document}

\section{\muRTL{} Proof Rules}

Having established a simulation relation,
the semantics of \muRTL{},
and a separation logic over pairs of memories,
we combine these components into a relational proof system.
To facilitate the verification of code transformations,
we establish a structured collection of relational proof rules
operating directly on the simulation judgment $\fsim{t}{j}{i}{s}{\Phi}$.
These rules are all proved sound with respect to the simulation relation
and the \muRTL{} semantics.

\subsection{Asynchronous Rules}

Our simulation relation decouples the execution of the target program
from that of the source program.
This separation yields asynchronous proof rules that advance one program
while keeping the configuration of the other fixed.

Target proof rules express the conditions under which the target program
can execute an instruction.
The rules governing target execution of loads, allocations, and operations
are displayed in \cref{fig:target-rules}.
The source stepping rules mirror the target rules,
operating on the source program counter,
source registers,
and source memory assertions.

\begin{figure}[htp]
  \centering
  \begin{gather*}
    \begin{prooftree}
      \hypo{
        \begin{matrix}
          \evalop{\textsf{op}}{\vec{v}} = \Some{v}
          \\[.2em]
          \pcatis{f}{pc}{\Iop{dst}{\textsf{op}}{\vec{r}}{pc'}}
        \end{matrix}
      }
      \hypo{
        \begin{matrix}
          \getreg{\vec{r}}{\rho} = \vec{v}
          \\[.2em]
          \fsim{(\sigma, \State{f}{pc'}{\setreg{dst}{v}{\rho}})}
          {j'}{i}{c}{\Phi}
        \end{matrix}
      }
      \infer2[\textsc{Target-Op}]{
        \fsim{(\sigma, \State{f}{pc}{\rho})}{j}{i}{c}{\Phi}
      }
    \end{prooftree}
    \\[1em]
    \begin{prooftree}
      \hypo{
        \begin{matrix}
          \pcatis{f}{pc}{\Iload{dst}{src}{pc'}}\qquad
          \getreg{src}{\rho} = \VPtr{l}\qquad
          \ptrt{l}{v}
          \\[.2em]
          \ptrt{l}{v} \wand \fsim{(\sigma, \State{f}{pc'}{\setreg{dst}{v}{\rho}})}
          {j'}{i}{c}{\Phi}
        \end{matrix}
      }
      \infer1[\textsc{Target-Load}]{
        \fsim{(\sigma, \State{f}{pc}{\rho})}{j}{i}{c}{\Phi}
      }
    \end{prooftree}
    \\[1em]
    \begin{prooftree}
      \hypo{
        \begin{matrix}
          \pcatis{f}{pc}{\Ialloc{dst}{pc'}}
          \\[.2em]
          \forall l\, v,\;
          \ptrt{l}{v} \wand
          \fsim{(\sigma, \State{f}{pc'}{\setreg{dst}{\VPtr{l}}{\rho}})}
          {j'}{i}{c}{\Phi}
        \end{matrix}
      }
      \infer1[\textsc{Target-Alloc}]{
        \fsim{(\sigma, \State{f}{pc}{\rho})}{j}{i}{c}{\Phi}
      }
    \end{prooftree}
  \end{gather*}
  \caption{
    Asynchronous rules for allocation
    and operation instructions in the target program.
  }
  \label{fig:target-rules}
\end{figure}

A key semantic difference between target and source execution appears
in the allocation rule \textsc{Target-Alloc} versus \textsc{Source-Alloc}.
In our operational semantics,
allocating a fresh memory cell non-deterministically initializes its content
with an arbitrary value.
Because the target is demonic,
the continuation in \textsc{Target-Alloc} must hold universally for
\emph{all} possible fresh locations $l$ and all initial values $v$.
In contrast,
because the source is angelic,
when applying \textsc{Source-Alloc},
the value of the allocated memory cell may be chosen existentially:
\begin{equation*}
  \begin{prooftree}
    \hypo{
      \begin{matrix}
        v \in \val \qquad \pcatis{f}{pc}{\Ialloc{dst}{pc'}}
        \\[.2em]
        \forall l,\;
        \ptrs{l}{v} \wand
        \fsim{c}{j}{i'}
        {(\sigma, \State{f}{pc'}{\setreg{dst}{\VPtr{l}}{\rho}})}
        {\Phi}
      \end{matrix}
    }
    \infer1[\textsc{Source-Alloc}]{
      \fsim{c}{j}{i}{(\sigma, \State{f}{pc}{\rho})}{\Phi}
    }
  \end{prooftree}
\end{equation*}

\subsection{Undefined Behavior Exploitation}

By the \textsc{FSourceStuck} rule, our free simulation is skewed.
When the source program reaches a stuck configuration,
the simulation property $\fsim{t}{j}{i}{s}{\Phi}$ holds directly.
This property provides a powerful proof technique:
whenever the source program is about to step,
we can safely assume that the configuration is not stuck.
Otherwise the proof obligation would be instantly discharged.
For instance,
the exploit rule \textsc{Source-Op-Exploit} allows the proof to assume
that the operation succeeds and produces a value $v$:
\begin{equation*}
  \begin{prooftree}
    \hypo{
      \begin{matrix}
        \getreg{\vec{r}}{\rho} = \vec{v} \qquad
        \pcatis{f}{pc}{\Iop{dst}{\textsf{op}}{\vec{r}}{pc'}}
        \\[.2em]
        \forall v,\;
        \purel{\evalop{\textsf{op}}{\vec{v}} = \Some{v}}
        \wand
        \fsim{c}{j}{i'}
        {(\sigma, \State{f}{pc'}{\setreg{dst}{v}{\rho}})}
        {\Phi}
      \end{matrix}
    }
    \infer1[\textsc{Source-Op-Exploit}]{
      \fsim{c}{j}{i}{(\sigma, \State{f}{pc}{\rho})}{\Phi}
    }
  \end{prooftree}
\end{equation*}

\subsection{Exploiting Memory Operations}

In separation logic,
memory resources cannot be obtained without explicit allocation.
However,
when programs manipulate values related through a memory injection
$\meminj{I}{E}$,
heap assertions can be extracted directly
from the global injection invariant.

To justify the soundness of the \textsc{Source-Load-Exploit} rule
presented in \cref{fig:source-load-exploit},
suppose the source program is about to execute a load on an address $v$
that is related to a target value $w$ through $\sameval{I}{w}{v}$.
The proof sketch proceeds in two steps:
\begin{enumerate}
  \item \textbf{Pointer validity:}
        By case analysis,
        if $v$ is not a valid pointer $\VPtr{l}$,
        the source becomes stuck
        and the proof is immediate via \textsc{FSourceStuck}.
        Otherwise,
        $v = \VPtr{l}$,
        and $\sameval{I}{w}{v}$ ensures that $w = \VPtr{l'}$
        with $(l', l) \in I$.
  \item \textbf{Active status:}
        Opening the memory injection $\meminj{I}{\emptyset}$
        via the extraction \cref{lem:inj-extraction}
        reveals that the pair $(l', l)$
        is either actively allocated or already freed.
        If the location were freed,
        attempting a load in the source would produce a stuck configuration.
        Thus,
        the freed case is instantly discharged,
        and the cell is guaranteed to be currently allocated.
\end{enumerate}

From a user perspective,
this rule exhibits an interesting asymmetry.
In standard separation logic,
and in our \textsc{Target-Load} rule,
executing a load demands an existing points-to assertion in the context.
In contrast,
\textsc{Source-Load-Exploit} does not consume any points-to assertion.
Instead,
it extracts and provides both points-to assertions
$\ptrt{l'}{a} \sep \ptrs{l}{b}$
from the global invariant $\meminj{I}{\emptyset}$,
adding $(l', l)$ to the exception set.
For clarity,
\cref{fig:source-load-exploit} displays a simplified rule
where the exception set is initially empty.
Similar exploitation principles apply to the remaining memory operations.

\begin{figure}[htp]
  \begin{equation*}
    \begin{prooftree}
      \hypo{
        \begin{matrix}
          \pcatis{f}{pc}{\Iload{dst}{src}{pc'}}
          \\[.2em]
          {
          \getreg{src}{\rho} = v
          \qquad
          \sameval{I}{w}{v}
          \qquad
          \meminj{I}{\emptyset}
          }
          \\[.2em]
          \left(
          \begin{array}{@{}l@{}}
            \forall l'\, l\, a\, b,\;
            v = \VPtr{l} \land
            w = \VPtr{l'} \implies {}
            \\
            \hspace{5em}
            \ptrt{l'}{a} \sep
            \ptrs{l}{b} \sep
            \purel{\sameval{I}{a}{b}} \sep
            \meminj{I}{\{ (l', l)\}} \wand {}
            \\
            \hspace{5em}
            \fsim{c}{j}{i'}
            {(\sigma, \State{f}{pc'}{\setreg{dst}{b}{\rho}})}{\Phi}
          \end{array}
          \right)
        \end{matrix}
      }
      \infer1[\textsc{Source-Load-Exploit}]{
        \fsim{c}{j}{i}{(\sigma, \State{f}{pc}{\rho})}{\Phi}
      }
    \end{prooftree}
  \end{equation*}
  \caption{Memory exploitation rule for source load instructions.}
  \label{fig:source-load-exploit}
\end{figure}

\subsection{Synchronous Rules}

While asynchronous and exploit rules provide flexibility for transformations
that reorder or eliminate instructions,
many compiler passes preserve the overall control flow of the program.
For such passes,
stepping target and source programs independently introduces redundant
bookkeeping.
By composing source exploit rules with target stepping rules,
we obtain synchronous rules that execute matching instructions
in lockstep.

This is notably the case when reasoning about synchronous function calls.
To ensure modularity,
we avoid explicitly stepping through the execution of the callees.
Instead,
we rely on the Hoare quadruples introduced in \cref{sec:rel-hoare}.
The synchronous call rule \textsc{Both-Call} enables replacing a pair of
matching function calls by their relational specification.
Its structure highlights the core benefits of our relational separation logic.

Suppose both programs are about to execute a call:
\begin{itemize}
  \item the target calls function $f_t'$ with argument values $\vec{v}_t$,
        where $\pcatis{f_t}{pc_t}{\Icall{dst_t}{f_t'}{\vec{a}_t}{pc_t'}}$
        and $\getreg{\vec{a}_t}{\rho_t} = \vec{v}_t$;
  \item the source calls function $f_s'$ with argument values $\vec{v}_s$,
        where $\pcatis{f_s}{pc_s}{\Icall{dst_s}{f_s'}{\vec{a}_s}{pc_s'}}$
        and $\getreg{\vec{a}_s}{\rho_s} = \vec{v}_s$;
  \item a relational specification $\hoareq{P}{f_t'}{f_s'}{Q}$
        has already been established for the called functions.
\end{itemize}

Under these conditions,
the \textsc{Both-Call} rule proceeds as follows:
\begin{enumerate}
  \item The caller must prove that current heap resources satisfy
        the relational precondition $P(\vec{v}_t, \vec{v}_s)$.
  \item Instead of stepping through the internal instructions of $f_t'$
        and $f_s'$,
        the rule bypasses the calls entirely.
  \item Execution resumes at the return points $pc_t'$ and $pc_s'$
        with the destination registers updated with return values,
        satisfying the postcondition $Q(v_t, v_s)$.
\end{enumerate}

Because Hoare quadruples are enclosed in the persistence modality,
they represent duplicable specifications
that can be verified once per function,
and reused across many call sites.
This provides horizontal modularity:
functions in a program can be compiled and verified independently,
and their correctness guarantees compose cleanly through \textsc{Both-Call}.

\end{document}

% LocalWords:  Coinductive coinductive asymmetricity arities
